\documentclass[12pt]{article}
\usepackage[utf8]{inputenc}
\usepackage[spanish]{babel}
\usepackage[a4paper, margin=2cm]{geometry}
\usepackage{tikz}
\usetikzlibrary{babel}
\usepackage[most]{tcolorbox}
\usepackage{amsmath}

% Usar fuente sin serifa para un aspecto más moderno
\renewcommand{\familydefault}{\sfdefault}

% Definición de colores
\definecolor{medblue}{RGB}{43, 108, 176}
\definecolor{medteal}{RGB}{49, 151, 149}
\definecolor{medlight}{RGB}{235, 248, 255}
\definecolor{meddark}{RGB}{26, 32, 44}
\definecolor{nerveyellow}{RGB}{245, 158, 11} % Amarillo/Naranja para la neurona
\definecolor{musclepink}{RGB}{244, 114, 182} % Rosa para el músculo

\begin{document}
\pagestyle{empty}

\begin{tcolorbox}[
    enhanced, 
    colback=medlight!30, 
    colframe=medblue, 
    title={\Large \textbf{Cinemática: Impulso Nervioso (MRU)}}, 
    halign title=center, 
    fonttitle=\bfseries, 
    arc=4mm, 
    boxrule=1.5pt,
    shadow={2mm}{-2mm}{0mm}{black!30!white}
]

    \vspace{0.2cm}
    \begin{tcolorbox}[colback=white, colframe=medteal, arc=2mm, boxrule=1.2pt, title=\textbf{Caso Clínico / Problema}]
        Un impulso nervioso se propaga a lo largo de un axón a una velocidad constante de \textbf{120 m/s}. Calcule el tiempo que tarda el impulso en recorrer una distancia de \textbf{0.8 m} desde la médula espinal hasta un músculo de la pierna.
    \end{tcolorbox}

    \vspace{0.4cm}
    
    \begin{center}
    \begin{tikzpicture}[scale=0.95]
        
        % Médula (izquierda)
        \filldraw[medblue!20, draw=medblue, thick, rounded corners=5pt] (-2.5, -1.5) rectangle (-0.5, 1.5);
        \node[font=\bfseries, text=meddark, align=center] at (-1.5, 0) {Médula\\Espinal};

        % Neurona (Soma)
        \filldraw[nerveyellow!50, draw=nerveyellow!90!black, thick] (0,0) circle (0.5cm);
        % Dendritas
        \foreach \a in {30, 90, 150, 210, 270, 330} {
            \draw[thick, nerveyellow!90!black] (0,0) -- (\a:0.8);
        }
        
        % Axón
        \filldraw[nerveyellow!30, draw=nerveyellow!90!black, thick] (0.4,-0.1) rectangle (6,0.1);
        
        % Vainas de mielina
        \foreach \x in {1, 2.5, 4} {
            \filldraw[medteal!30, draw=medteal, thick, rounded corners=2pt] (\x, -0.25) rectangle (\x+1, 0.25);
        }

        % Músculo (derecha)
        \filldraw[musclepink!50, draw=musclepink!90!black, thick, rounded corners=3pt] (6, -1.2) rectangle (7.5, 1.2);
        \draw[thick, musclepink!90!black] (6.3, -1) -- (6.3, 1);
        \draw[thick, musclepink!90!black] (6.7, -1) -- (6.7, 1);
        \draw[thick, musclepink!90!black] (7.1, -1) -- (7.1, 1);
        \node[font=\bfseries, text=meddark, align=center] at (6.75, -1.6) {Músculo};

        % Impulso nervioso (Vector de Velocidad)
        \draw[->, ultra thick, red!80!black] (1.5, 0.8) -- (3.5, 0.8) node[midway, above, font=\bfseries] {Impulso: $v = 120 \text{ m/s}$};

        % Cota de Distancia
        \draw[<->, thick, meddark] (0, -1.2) -- (6, -1.2) node[midway, below, font=\bfseries] {$d = 0.8 \text{ m}$};

    \end{tikzpicture}
    \end{center}

    \vspace{0.4cm}
    \begin{tcolorbox}[colback=medteal!10, colframe=medteal, arc=2mm, boxrule=1.2pt, title=\textbf{Desarrollo Matemático y Solución}]
        Dado que el impulso viaja a velocidad constante, aplicamos la ecuación del Movimiento Rectilíneo Uniforme (MRU) $v = \frac{d}{t}$. Despejando el tiempo $t$:
        \begin{align*}
            t &= \frac{d}{v} = \frac{0.8 \text{ m}}{120 \text{ m/s}} \\[1ex]
            t &\approx \mathbf{0.00667 \text{ s}} \quad \text{o} \quad \mathbf{6.67 \text{ ms}}
        \end{align*}
        \textbf{Resultado:} El impulso nervioso tarda aproximadamente \textbf{6.67 milisegundos} en llegar al músculo.
    \end{tcolorbox}
    
\end{tcolorbox}

\end{document}